% Chapter 14: Acute and Obtuse Triangle Cases
\chapter{Acute and Obtuse Triangle Cases}
\label{ch:acute-obtuse}

\section{Classification Criteria}
A triangle $\tri ABC$ is classified by its largest angle:
\begin{itemize}
    \item \textbf{Acute:} All angles $< \pi/2$, equivalently $a^2 + b^2 > c^2$, $b^2 + c^2 > a^2$, $a^2 + c^2 > b^2$.
    \item \textbf{Right:} One angle $= \pi/2$ (Chapter~\ref{ch:right}).
    \item \textbf{Obtuse:} One angle $> \pi/2$, equivalently one of the squared-side inequalities is reversed.
\end{itemize}

\section{Acute Triangle Constraints}
For an acute triangle, all three cosine inequalities hold:
\begin{equation}
\cos\hat{A} > 0, \quad \cos\hat{B} > 0, \quad \cos\hat{C} > 0.
\end{equation}
Using $\cos\hat{A} = (b^2 + c^2 - a^2)/(2bc)$, the acute conditions translate to:
\begin{equation}
a^2 < b^2 + c^2, \quad b^2 < a^2 + c^2, \quad c^2 < a^2 + b^2.
\end{equation}
These three strict inequalities carve out a \emph{proper open cone} in $(a, b, c)$-space. The intersection of the solution manifold $\mathcal{M}$ with this cone gives the set of acute triangles compatible with the given vertex distances.

\section{Obtuse Triangle Constraints}
For an obtuse triangle with $\hat{C} > \pi/2$ (WLOG, the largest angle is at $C$):
\begin{equation}
c^2 > a^2 + b^2.
\end{equation}
This reverses one inequality, placing the solution in a different region of $(a, b, c)$-space.

\section{Impact on the Subtended Angles}

The type of triangle affects the range of subtended angles at $P$:
\begin{proposition}
For an acute triangle $\tri ABC$ with interior point $P$:
\begin{equation}
\frac{\pi}{3} < \alpha, \beta, \gamma < \pi \quad \text{(with } \alpha + \beta + \gamma = 2\pi\text{)}.
\end{equation}
For an obtuse triangle, the subtended angle opposite the obtuse vertex can be smaller than $\pi/3$.
\end{proposition}

\section{Solution Manifold Decomposition}
The solution manifold $\mathcal{M}$ decomposes into three regions based on triangle type:
\begin{equation}
\mathcal{M} = \mathcal{M}_{\text{acute}} \sqcup \mathcal{M}_{\text{right}} \sqcup \mathcal{M}_{\text{obtuse}},
\end{equation}
where $\mathcal{M}_{\text{right}}$ consists of isolated points (codimension 1 in $\mathcal{M}$) and the other two are open arcs.

\section{Numerical Comparison}
The computational engines (Chapter~\ref{ch:numerical}) classify each solution found on $\mathcal{M}$ by type and compute type-specific statistics:
\begin{itemize}
    \item Fraction of the manifold arc that corresponds to each type.
    \item Perimeter range within each type.
    \item Location of the type transitions (right triangle boundaries).
\end{itemize}
